To illustrate how measured performance can inform a resource-constrained choice, an offline screening example uses a prespecified eight-candidate pool, fixed before the paired-subset completion. A candidate $c$ specifies the actuated node and its observation subset. The pool comprises Nodes 2 and 3, each with local two-dimensional observations, the original four-dimensional subset ($\mathcal O_{12}$ at Node 2 and $\mathcal O_{13}$ at Node 3), $\mathcal O_{23}$, and full six-dimensional observations. Each candidate retains all three training runs; configuration selection therefore accounts for variability across the retained runs. The allowed actuator set may be $\{2,3\}$, $\{2\}$, or $\{3\}$. The subsequently evaluated Node 2/$\mathcal O_{13}$ and Node 3/$\mathcal O_{12}$ configurations are not included among these eight candidates; the complete paired-subset comparison and this limited selection example answer different questions.

For a trajectory, let $z$ be the sample mean of $\|\mathbf e(t)\|_1$ over the final interval $16\leq t\leq20$. A trajectory passes when $z\leq\varepsilon$ and no state-protection boundary is reached. For training run $r$, its validation pass fraction is
\begin{equation}
 \widehat p_{c,r}(\varepsilon)
 =\frac{1}{n_v}\sum_{\ell=1}^{n_v}
 \mathbf 1\{z_{c,r,\ell}\leq\varepsilon,\ g_{c,r,\ell}=0\},
 \label{eq:configuration_pass}
\end{equation}
where $g$ counts state-protection boundary hits. A candidate is feasible only if $\min_r\widehat p_{c,r}\geq0.90$. Among feasible candidates with allowed actuation, the rule first minimizes observation dimension and then validation mean applied energy $E_u$. Remaining ties are resolved by a fixed candidate ordering. If none is feasible, no configuration is selected. Thus, fewer observed variables take priority over lower energy once the accuracy requirement is met.

The rule uses 20 new common validation initial conditions and 40 separate common test initial conditions drawn from the original initial-state distribution. The illustrative tolerances $\varepsilon\in\{0.01,0.05,0.10\}$ were informed by the preceding experiments and fixed before this validation. They serve as illustrative design targets for the screening example. Choices were recorded before evaluating the selected configurations on the test set, with no test-based reselection. Dynamics, nominal parameters, action bounds, and policy parameters were unchanged.

\begin{table*}[pos=!t,width=\textwidth]
\centering
\caption{Configurations selected using validation data and evaluated on separate test initial conditions. Test pass fractions are the minimum across three training runs, each evaluated on the same 40 test initial conditions. Error and energy average initial conditions and runs. A dash denotes no selection and hence no corresponding test; grouped tolerances have identical outcomes.}
\label{tab:configuration_selection}
\small
\begin{tabular}{cccccc}
\toprule
Allowed nodes & $\varepsilon$ & Selected node / observations & Test pass & Test $z$ & Test $E_u$\\
\midrule
$\{2,3\}$ & 0.01 & -- & -- & -- & --\\
$\{2,3\}$ & 0.05, 0.10 & $3\,/\,\mathcal O_{13}$ & 100\% & 0.0086 & 5.03\\
$\{2\}$ & 0.01, 0.05 & -- & -- & -- & --\\
$\{2\}$ & 0.10 & $2\,/\,\mathcal O_{23}$ & 100\% & 0.0282 & 8.11\\
$\{3\}$ & 0.01 & -- & -- & -- & --\\
$\{3\}$ & 0.05, 0.10 & $3\,/\,\mathcal O_{13}$ & 100\% & 0.0086 & 5.03\\
\bottomrule
\end{tabular}
\end{table*}

Table~\ref{tab:configuration_selection} translates the observation comparison into conditional design choices. With either node available, the original four-dimensional Node 3 configuration is selected at tolerances 0.05 and 0.10. When actuation is restricted to Node 2, the alternative four-dimensional subset is selected at 0.10. Both selected configurations pass every test trajectory at their respective selected tolerances. No candidate in the prespecified pool meets the validation criterion at 0.01. No state-protection boundary is reached in the 480 validation or 240 test trajectories.

The per-run criterion prevents aggregate averages from masking inconsistent training outcomes. At tolerance 0.05, the alternative Node 2 subset has validation pass fractions of 100\%, 75\%, and 100\%; it is rejected despite an aggregate fraction above 90\%. The local two-dimensional Node 3 policies likewise yield 100\%, 0\%, and 100\% at both 0.05 and 0.10. The resource priority affects the choice: at Node 2, the selected four-dimensional configuration has validation mean energy 11.76 versus 10.93 for six dimensions. The preference must therefore reflect the intended resource priority.

The procedure selects among the eight candidates according to nominal accuracy, observation dimension, and energy, then checks the choice on separate test initial conditions. Pass fractions describe the tested policies and sampled initial conditions; they do not establish the same pass rates for all possible initial conditions or newly trained policies. The tail-mean criterion measures average late-stage accuracy, whereas pointwise error limits and application-specific recovery requirements call for their own acceptance criteria. Shared initial conditions provide paired comparisons across policies.

To examine whether the choice depends on the selection priorities, an additional analysis uses the same eight-candidate validation data, varying the minimum per-run pass fraction over 80\%, 90\%, and 100\%, and comparing dimension-first with energy-first ranking. Changing the pass-fraction requirement does not change any choice at the tested error tolerances. The ranking preference matters when only Node 2 is allowed at $\varepsilon=0.10$: dimension-first selects $\mathcal O_{23}$, whereas energy-first selects full observations. All other choices are unchanged. This analysis uses validation data only and does not alter the configurations already chosen for independent testing. The subsequently completed paired subsets are not retroactively inserted into the original candidate set or the previously selected configurations.

We then examine whether the nominal choices retain their requested accuracy when response parameters change. The original selected Node 2/$\mathcal O_{23}$ and Node 3/$\mathcal O_{13}$ policies are evaluated without changing their parameters at their respective tolerances 0.10 and 0.05. For 40 new common initial conditions, response weights are set to $W_r=W\odot(1+\eta U)$, where each entry of $U$ is sampled uniformly from $[-1,1]$ and $\eta\in\{0,0.01,0.03,0.05\}$. The drive system is unchanged. Each initial condition uses the same mismatch direction across amplitudes and policies; an independent random stream preserves identical initial states. Zero weights remain zero. No retraining or configuration reselection is performed.

\begin{table*}[pos=!t,width=\textwidth]
\centering
\caption{Performance of the selected configurations under response-weight mismatch, without retraining or reselection. Pass fractions are the minimum across three training runs, each evaluated on the same 40 new initial conditions. Mean tail error $z$ and energy $E_u$ average runs and initial conditions. Pass fractions track each configuration against its own previously fixed tolerance.}
\label{tab:configuration_transfer}
\small
\begin{tabular}{cccccc}
\toprule
Configuration & $\varepsilon$ & $\eta$ & Minimum pass & Mean $z$ & Mean $E_u$\\
\midrule
$2/\mathcal O_{23}$ & 0.10 & 0\% & 100.0\% & 0.0266 & 10.08\\
$2/\mathcal O_{23}$ & 0.10 & 1\% & 97.5\% & 0.0402 & 10.11\\
$2/\mathcal O_{23}$ & 0.10 & 3\% & 57.5\% & 0.0942 & 10.20\\
$2/\mathcal O_{23}$ & 0.10 & 5\% & 35.0\% & 0.1571 & 10.39\\
$3/\mathcal O_{13}$ & 0.05 & 0\% & 100.0\% & 0.0084 & 5.84\\
$3/\mathcal O_{13}$ & 0.05 & 1\% & 97.5\% & 0.0181 & 5.85\\
$3/\mathcal O_{13}$ & 0.05 & 3\% & 75.0\% & 0.0455 & 5.91\\
$3/\mathcal O_{13}$ & 0.05 & 5\% & 42.5\% & 0.0743 & 6.01\\
\bottomrule
\end{tabular}
\end{table*}

Table~\ref{tab:configuration_transfer} shows that both configurations retain a minimum pass fraction of 97.5\% at 1\% relative mismatch, but fall below the original 90\% requirement at 3\% and 5\%. No state-protection boundary is reached in these 960 trajectories. The drop in pass fraction identifies the sensitivity of the selected configurations to larger response-weight deviations and motivates incorporating the intended parameter range into configuration validation.

